% ===========================================================================
\section{Inventory construction}
\label{sec:methods}
% ===========================================================================

\subsection{Segment ledger and practice selection}

For segment $s$ and gas $g$, emissions are
\begin{equation}
E_{s,g}=A_s EF_{s,g}.
\end{equation}
Each ledger row carries the IPCC code, activity and unit, source and imputation status, factor,
uncertainty bounds and structural configuration. The accounting engine consumes this generic
schema, allowing facility records to replace aggregate proxies without changing the downstream
calculations.

The \emph{reference configuration} uses the higher/default 2019 practice factors where national
LDAR, pipeline-material or compressor-seal evidence is unavailable. Lower-practice factors and
throughput-based transmission and distribution are evaluated as structural alternatives.
Produced gas provides the processing proxy; sour-gas processing uses the 2019 default
\SI{32}{\percent} share. Both assumptions can be varied or excluded independently.

Post-meter appliances use one appliance per domestic meter and the Refinement's warm-country
fallback. The industrial/power factor is \SI{6.1}{\tonne} CH$_4$ per million m$^3$, derived in
the 2019 Refinement from a Revised 1996 value of \SI{175000}{\kilogram\per\peta\joule} and a
\SI{35}{\mega\joule\per\cubic\metre} gas conversion. Its older empirical basis motivates full,
half-factor and exclusion configurations. The complete factor lineage and uncertainty register
are given in the supplement.

\subsection{Method-resolution experiment}

Five estimates separate factor changes from boundary expansion:
\begin{enumerate}[leftmargin=*,itemsep=2pt]
\item the legacy 2006 production-only calculation for FY2019--20;
\item the same method on FY2023--24 production;
\item 2019 factors on identical FY2023--24 oil and gas production;
\item the 2019 upstream account after adding gathering and processing; and
\item the public-data 2019 ledger, including T\&D, LNG, refining and post-meter sources.
\end{enumerate}
The common-support comparison isolates the factor-framework effect; the later steps measure the
effect of representing additional physical segments.

\subsection{Monthly timing, upstream only}

The national release resolves this sector into fiscal months with process-specific support. Validated PBS provincial natural-gas and crude-oil production can support corresponding upstream production, gathering and processing timing as a throughput proxy. Twelve rows alone do not establish a complete national observation: source province coverage, units and compatible monthly or annual controls are checked against the final exported panel. Repaired extraction and worksheet selection are reflected in the current profiles. Where the source remains incomplete, an explicit fallback replaces the broken profile and carries weaker evidence; it is not labelled national observed activity.

Distribution, transmission, post-meter leakage and LNG follow deliveries or send-out, which are different physical drivers from field production. In the absence of appropriate observed monthly controls, those streams remain days-weighted or use a documented supported proxy. Refinery intake, wellhead production, flaring and venting retain distinct timing targets. The reference-year allocation does not imply observed changing geography across the extended annual series.

\subsection{Air-pollutant source boundaries}
The common source account supports admitted NMVOC and other pollutant factors for the specified oil/gas processes and flaring. Flaring combustion and venting are distinct, and factor coverage does not establish that every gas balance discrepancy is methane leakage. NH$_3$ and unresolved venting, post-meter or other pollutant categories remain notation gaps where no supported factor exists. The air ensemble reuses native fugitive activity realisations; an industrial gas-combustion draw deducts the same post-meter leakage quantity before applying a combustion factor.

\subsection{Process-matched allocation}

Production and gathering follow provincial oil and gas production. Transmission and
distribution follow measured pipeline length, post-meter industrial emissions follow gas
consumption, refineries follow facility output, and LNG terminals remain assigned to Sindh. Each
allocation reconciles exactly to its national control. Facility and network geometry is
insufficient for a defensible district grid, so the spatial product stops at process-matched
province and point layers.

The implementation tests unit exponents, source-row and segment closure, activity-year anchors,
pipeline and gas-flow identities, structural-model enumeration, uncertainty ordering,
observation coverage, official-inventory reproduction, provincial reconciliation and
manuscript-output agreement. Detailed assertions and generated products are documented in the
supplement and reproducibility appendix.
